
\title{DHC: Dynamic Heterogeneous Context Learning with Relation Memory for Temporal Link Prediction}
\titlerunning{DHC for Temporal Link Prediction}

\author{Xuejie Xu\inst{1}\orcidID{0000-0003-4890-2110} 
\and Wenjin Zhang\inst{1}\orcidID{0009-0007-8375-9280}  
\and \rev{Liangzhe Han}\inst{1}
\and Jian Huang\inst{1}\orcidID{0000-0002-8223-5705} 
\and Leilei Sun\inst{1}\orcidID{0000-0002-0157-1716}}
\authorrunning{X. Xu et al.}
\institute{Beihang University, Beijing 100191, China
\email{\{x\_xuejie,liangzhehan,zhangwenjin,hj,leileisun\}@buaa.edu.cn}}
\maketitle

\begin{abstract}
Dynamic heterogeneous graphs are commonly found in academic networks, patent systems, and software ecosystems. 
They contain timestamped interactions, typed relations, and text-derived node attributes. 
These signals are complementary, but existing temporal graph models still encode relation types as static edge features. They therefore have limited ability to track relation semantics as the graph evolves. 
We propose DHC, a dynamic heterogeneous context-learning framework for temporal link prediction. 
DHC combines text-aware node initialization, relation-time message generation, node-relation memory, recent-neighbor temporal subgraph encoding, and a sequential event-memory branch.
The relation-aware context branch captures heterogeneous local evidence, while the event-memory branch preserves compact chronological interaction history. 
Then a gate fuses the two representations for link scoring. 
\rev{Empirical evaluations on AMiner-Paper, DBLP, Patent-US, and THGL-Software show that DHC is competitive under chronological evaluation; three datasets additionally use history-level cold-start views, while the reported THGL-Software ``inductive'' view is a disclosed temporal fallback.}
The results indicate that explicitly modeling dynamic relation semantics and local temporal context are useful for scholarly, patent, and software network mining.

\keywords{Dynamic Heterogeneous Graph \and Temporal Link Prediction \and Relation Memory \and Academic Data Mining}
\end{abstract}
